\vspace{-30pt}
\chapter{Project Objectives}
\label{ch:objectives}


The primary objective of the BELLONA project is to design an unmanned aerial system capable of locating, intercepting, and safely recovering free-floating uncooperative balloons in a controlled manner. This bridges the capability gap between high-cost military intervention and passive monitoring for uncooperative balloons. This project focuses on developing a solution that is tailored for civilian purposes, such that safety and cost-per-sortie are dominant success factors.

\vspace{-10pt}

\subsubsection{Strategic Goals}
Beyond specific technical thresholds, the project is guided by three strategic pillars:

\begin{itemize}
    \item Safety-First Interception: To move beyond destructive 
    solutions. The objective is to prove the feasibility of a recovery-based or controlled-deflation approach that minimizes risk to third parties on the ground.
    \item Economic Accessibility: To lower the barrier for airport authorities to manage their own airspace security. This involves designing a system with a low enough lifecycle cost to be kept on-site for immediate deployment.
    \item Engineering Rigor and Synthesis: A core objective is rigorous engineering reasoning and explainability of all decisions.
    
\end{itemize}
\vspace{-10pt}

\subsubsection{Definition of Scope}
To ensure the project remains manageable within the DSE framework, the following boundaries define the objective scope:
\begin{itemize}
    \item Focus: The design effort is focused on the entire system consisting of the interceptor drone, the capture mechanism and the ground control station.
    \item Phase: The objective is to reach a high-fidelity conceptual design. 
    
    Physical prototyping is considered secondary to the analytical validation of the system's performance.
\end{itemize}
\vspace{-10pt}
\subsubsection{Alignment with Stakeholder Needs}
The project objectives are derived directly from the needs of airport operators, national security agencies, and meteorological organizations. While specific technical requirements are detailed in Chapter \ref{ch:requirements}, the overarching objective is to provide these stakeholders with a "Low-Collateral Damage" alternative to existing counter-UAS technologies.